\chapter{Исследовательская программа: граф, сигнал, метрики}

\section{Почему серия вообще начинается с эмбриона}

В архитектуре RIN графовая память не возникает как чистая абстракция. Она
должна рано или поздно стать рабочим субстратом, по которому проходит сигнал,
куда коммитятся новые следы и из которого извлекается контекст. Поэтому уже
начиная со второго эксперимента задача формулируется двойственно:

\begin{enumerate}[leftmargin=2em]
\item нужен \textbf{эмбриональный граф}, достаточно реальный, чтобы на нем
вылезали содержательные дефекты;
\item нужен \textbf{контролируемый стенд}, достаточно малый и читаемый, чтобы на
нем можно было измерять не только факт прохождения сигнала, но и качество
retrieval.
\end{enumerate}

Отсюда вся дальнейшая архитектура серии:

\begin{itemize}[leftmargin=2em]
\item реальные источники -- для столкновения с настоящей структурной грязью;
\item золотые синтетические миры -- для локализации конкретных механизмов;
\item перенос с золотого мира на эмбрион -- как проверка того, переживает ли
формула выход из лабораторной колбы.
\end{itemize}

\section{Как выглядит базовый графовый объект}

Во всех сериях используется одна и та же базовая идея памяти как направленного
взвешенного графа. Узлы представляют концепты, события, состояния, а позднее --
эпизодические или служебные элементы. Ребра несут тип и вес.

Минимально полезные поля узла:

\begin{itemize}[leftmargin=2em]
\item текстовое содержимое;
\item тип узла (\texttt{concept}, \texttt{episode}, \texttt{claim},
\texttt{relationship\_state} и т.\,д.);
\item embedding в общем векторном пространстве;
\item метаданные происхождения и эпистемического статуса.
\end{itemize}

Минимально полезные поля ребра:

\begin{itemize}[leftmargin=2em]
\item источник и приемник;
\item тип отношения (\texttt{IsA}, \texttt{RelatedTo}, \texttt{UsedFor},
\texttt{Causes}, \texttt{HasProperty} и пр.);
\item вес проводимости.
\end{itemize}

Уже на этом уровне скрыта важная методологическая ловушка. Один и тот же граф
можно рассматривать по меньшей мере в трех разных ролях:

\begin{enumerate}[leftmargin=2em]
\item как \textbf{хранилище фактов и эпизодов};
\item как \textbf{среду распространения сигнала};
\item как \textbf{объект переписывания и консолидации}.
\end{enumerate}

Серия \texttt{exp2--exp9a} почти целиком посвящена второй роли, но время от
времени касается первой и третьей. Из-за этого полезно заранее различать
операторы чтения и записи.

\section{Как выглядит retrieval в базовой версии}

Пусть задан набор стартовых якорей. Им присваивается начальная амплитуда, после
чего по графу запускается затухающее распространение. В простейшей форме можно
думать о шаге диффузии как о правиле

\[
a_{t+1}(v)=\lambda \sum_{u \to v} a_t(u)\, P(u \to v),
\]

где:

\begin{itemize}[leftmargin=2em]
\item $a_t(v)$ -- уровень активации узла $v$ на шаге $t$;
\item $\lambda$ -- коэффициент затухания;
\item $P(u \to v)$ -- нормированная доля потока, идущая из $u$ в $v$.
\end{itemize}

В реальных ноутбуках используется не буквальная марковская цепь с бесконечным
пределом, а конечный процесс с отсечением слабых хвостов. Практически это
значит следующее:

\begin{enumerate}[leftmargin=2em]
\item в стартовые узлы впрыскивается фиксированный заряд;
\item на каждом шаге активны только узлы выше минимального порога;
\item выходящий поток делится между исходящими ребрами пропорционально их весу;
\item процесс идет ограниченное число шагов или до опустошения фронта.
\end{enumerate}

Из этого уже видно, почему серия так быстро вышла к проблеме хабов. Узел может
оказаться ярким по двум разным причинам:

\begin{enumerate}[leftmargin=2em]
\item он действительно хорошо объясняется данным запросом;
\item в него удобно сходится много путей, даже если по смыслу он слишком общий.
\end{enumerate}

\section{Три уровня retrieval-скоринга}

По ходу серии постепенно кристаллизуется удобное разложение retrieval на три
содержательно разных уровня.

\subsection{Локальный транспорт}

Это вопрос о том, \emph{дошел ли сигнал вообще} и как он растекся по имеющимся
путям. Здесь живут:

\begin{itemize}[leftmargin=2em]
\item исходные якоря;
\item затухание;
\item веса ребер;
\item типизация проводимости;
\item возможные дополнительные каналы вроде semantic edges.
\end{itemize}

\subsection{Локальное удержание}

Даже если два узла получили одинаковый приток, не обязано быть так, что они
одинаково \emph{удерживают} этот приток как локально значимую активацию. Отсюда
рождаются знаменатели семейства $R(B)$:

\[
\retained(B)=\frac{\incoming(B)}{R(B)}.
\]

На этом уровне обсуждаются:

\begin{itemize}[leftmargin=2em]
\item степень узла;
\item взвешенная исходящая сила;
\item активный fan-in;
\item локальная capacity как ``эффективный объем представления''.
\end{itemize}

\subsection{Глобальная структурная слава}

Наконец, даже хороший локальный отклик может оказаться слишком высоким просто
потому, что узел глобально знаменит. Тогда появляется еще одна поправка:

\[
\score(v\mid q)=\frac{\text{LocalResponse}(v\mid q)}{\text{Prior}(v)^\eta}.
\]

Именно здесь возникают in-strength, out-strength, PageRank и их сравнение.

\begin{takeaway}
Одно из важнейших достижений серии состоит в том, что эти три уровня перестали
смешиваться. В начале работ любой неуспех мог быть списан на ``плохие
эмбеддинги'' или ``неудачную формулу''. К середине серии стало ясно, что
транспорт, удержание и глобальная слава -- разные механизмы, и провал любого из
них не лечится простой правкой соседнего.
\end{takeaway}

\section{Роль embedding-пространства}

Embedding в этой серии играет несколько ролей сразу.

\begin{enumerate}[leftmargin=2em]
\item Он нужен для anchoring: чтобы query-слова можно было связать с узлами
графа.
\item Он нужен для семантического gate: чтобы весить релевантность узла текущему
запросу.
\item Он нужен для additive bypass в \texttt{8a}: similarity начинает напрямую
входить в итоговый score.
\item Он нужен для semantic-edge construction в \texttt{9a}: ближайшие соседи в
embedding-пространстве становятся кандидатами на дополнительные ребра.
\end{enumerate}

Это делает особенно опасной ситуацию, когда embedding-пространство само по себе
дефектно. Если косинус живет в узком конусе, generic nodes могут становиться
псевдохабами еще до того, как сигнал вообще вошел в графовую динамику. Поэтому
ветка \texttt{8b} возникает не как побочная эстетика, а как попытка убрать
важный конфаундер: не следует ли приписываемый графу дефект на самом деле
свойствам самой similarity-геометрии.

\section{Какими метриками пользовалась серия}

Ранняя часть серии опиралась на простые и очень понятные критерии:

\begin{itemize}[leftmargin=2em]
\item правильный узел должен быть выше отвлекающего;
\item broad hub не должен побеждать специфический ответ, если запрос не о нем;
\item наоборот, legitimate hub не должен быть задавлен там, где запрос именно о
нем.
\end{itemize}

Позднее появляются более тонкие метрики:

\begin{itemize}[leftmargin=2em]
\item число пройденных pairwise inequality;
\item число решенных запросов;
\item margin между лучшим правильным и лучшим distractor;
\item mean rank и MRR;
\item manifold-aware success вместо проверки одного канонического узла;
\item специальные danger-checks: hub reflood и RAG collapse.
\end{itemize}

Сюда стоит добавить два методологических принципа, выработанных по ходу серии.

\begin{methodbox}
\textbf{Принцип 1.} Нельзя ранжировать методы по ``суммарному margin'', если они
по-разному масштабируют одну и ту же базовую картину. Иначе получится ложная
победа просто более агрессивной нормировки.

\textbf{Принцип 2.} Как только граф начинает выдавать не один канонический узел,
а соседний семантический manifold, оценка должна это учитывать. Иначе мы будем
объявлять провалом ситуацию, в которой память по существу попала в правильный
район, но не в ту формулировку вершины, которую человек заранее выписал в
таблицу ответов.
\end{methodbox}

\section{Почему понадобились золотые миры}

Реальный граф нужен для реализма, но он почти бесполезен как чистый измеритель
retrieval-качества, если там нет ground truth. Поэтому серия несколько раз
возвращается к hand-authored worlds. Это не ``игра в песочнице'' в плохом
смысле, а нормальный исследовательский ход:

\begin{enumerate}[leftmargin=2em]
\item на реальном графе вылавливаются инфраструктурные и содержательные дефекты;
\item на золотом мире локализуется механизм;
\item затем гипотеза переносится обратно на более реальный срез.
\end{enumerate}

Слабость золотого мира известна заранее: он зависит от авторского конструирования.
Его сила тоже известна заранее: только он позволяет быстро различить, является
ли неуспех следствием плохой формулы, плохого графа или грязной постановки.

\section{Как устроен дальнейший рассказ}

Оставшиеся главы идут не просто по номерам, а по логическим сериям.

\begin{itemize}[leftmargin=2em]
\item Главы \texttt{exp2--exp5} строят базовые объекты: эмбриональный граф,
золотой мир, семантический gate и write-path.
\item Глава о \texttt{6x} разбирает локальное удержание и глобальные priors.
\item Глава о \texttt{7x} показывает, как эти идеи переносятся на реальный
эмбрион и где именно перенос начинает ломаться.
\item Глава о \texttt{8a--9a} фиксирует смещение от ``оптимизации знаменателя'' к
проблеме транспортной архитектуры и реконструкции путей.
\item Заключительная глава собирает общую карту того, что серия действительно
доказала и где ее естественная развилка дальше.
\end{itemize}
